\cleardoublepage

\newrefsection

\chapter{外文翻译}

\begin{center}
{\Large\bfseries 以太坊区块链金融机器人检测技术研究}
\end{center}

\subsection*{摘要}

机器人在分布式账本技术（DLT）中的集成促进了效率和自动化。然而，它们的使用也与掠夺性交易和市场操纵有关，并可能对系统完整性构成威胁。因此，了解机器人在DLT中的部署程度至关重要；尽管如此，当前的检测系统主要是基于规则的，缺乏灵活性。在本研究中，我们提出了一种利用机器学习检测以太坊平台上金融机器人的新方法。首先，我们系统化现有的科学文献并收集轶事证据，以建立一个金融机器人的分类法，包含7个大类和24个子类。接下来，我们创建了一个包含133个人类地址和137个机器人地址的真实数据集。第三，我们采用无监督和有监督的机器学习算法来检测部署在以太坊上的机器人。性能最高的聚类算法是高斯混合模型，平均聚类纯度为82.6\%，而性能最高的二分类模型是随机森林，准确率为83\%。我们基于机器学习的检测机制通过提供对当前机器人景观的额外见解，有助于理解以太坊生态系统动态。

\subsection*{关键词}

机器人，以太坊，区块链，DeFi，机器学习

\subsection*{1 引言}

金融活动的算法自动化是去中心化金融（DeFi）所利用的最根本的创新之一，并且对于在此框架内实现效率至关重要[3]。在像以太坊这样的分布式账本技术（DLT）中，机器人促进了被称为智能合约的软件程序的自动执行，这些程序封装了诸如传统金融操作等确定性函数的逻辑。机器人通常管理外部拥有账户（EOA），通常被称为钱包，这些钱包由一个十六进制地址标识。由于以经济利益、节省时间或提高容错性的形式自动化EOA存在显著激励，机器人已成为DLT上有影响力的参与者，并且其相关性随时间推移变得更加突出。例如，机器人越来越多地被利用，迄今为止已从以太坊生态系统中提取超过10亿美元的利润[2, 16, 29]。

机器人的部署对基于DLT的生态系统可能是有益的。机器人可用于提供关键基础设施，例如，通过自动管理资产，实现中心化交易所与区块链之间的接口[33]。此外，像roll-ups这样的第二层扩展解决方案会自动将交易信息部署到以太坊上以实现持久性[38]。机器人也通过套利为市场稳定做出贡献[42]。

然而，机器人也被用于掠夺性交易和市场操纵[29]，对不警惕的用户构成财务威胁[43]，并对网络的完整性构成潜在的系统性威胁[5]。在利润优化型机器人的情况下，对生态系统的副作用可能包括利用与智能合约交互的用户[39]、利用智能合约[47]、通过增加的流量抬高Gas价格，以及网络的不稳定[13]。

虽然了解DLT中机器人采用的规模对于评估对人类用户的风险、确保他们的保护以及为政策提供信息至关重要，但很难估计这种现象的程度，因为迄今为止还没有专门的机器人检测系统。现有的软件，如MEV-inspect和Eigenphit，专注于通常由机器人执行的利润优化交易，可能被重新用于机器人检测。然而，它们是基于规则的，并非为捕捉不断演变的机器人行为而构建。因此，设计更灵活、以数据驱动方式运行的机器人检测方法非常重要。

本研究提供了一种基于机器学习的方法来检测部署在以太坊上的金融机器人，并旨在识别关键的机器人预测因子，以消除目前对特定机制和基于规则系统的普遍依赖。图1展示了我们工作的流程概览。首先，我们用轶事证据补充现有文献，提出以太坊上金融机器人的分类法。其次，我们处理公开可用的以太坊区块链数据以及来自Etherface.io和Etherscan.io网站的额外信息，并基于三位独立评估者的评估，通过注释270个EOA创建了一个真实数据集。然后，我们提取特征以准确表示EOA行为，用于使用机器学习算法测试机器人检测。最后，我们评估这些模型的有效性。我们利用无监督和有监督技术分别将EOA分组和分类为“机器人”或“人类”。此外，我们关注资产积累型机器人的一个特定子集，并在多类分类背景下衡量算法性能。本研究的研究问题和相关贡献可总结如下：

RQ1 在像以太坊这样的分布式账本技术上活跃着哪些类型的机器人？我们为以太坊上的金融机器人提出了一个分类法，包含7个机器人大类和24个子类。此外，我们创建了一个包含133个人类地址和137个机器人地址的真实数据集。

RQ2 我们能在多大程度上自动检测机器人？我们部署了有监督和无监督的机器学习模型并衡量了它们的性能。性能最高的聚类算法是高斯混合模型，在平衡的二类数据集上评估的平均聚类纯度为82.6\%。性能最高的二分类模型是随机森林，在相同数据集上的准确率为83\%。

RQ3 哪些特征对高性能模型最具信息量？使用可解释人工智能技术，我们发现基于传出交易的时间、频率、Gas价格和Gas限制的特征最具影响力。

我们相信，我们的工作通过揭示现有机器人景观的更多细节并提出一种新颖的基于机器学习的检测机制，有助于更好地理解以太坊生态系统动态。由于这些机器人可以显著影响市场动态、流动性和整体网络安全，监控它们的普遍性和影响非常重要。

本研究中使用的所有数据和源代码均为公开可用的，以便于重现结果。

\subsection*{2 背景}

本节提供了金融机器人所使用的技术的背景，定义了区块链背景下机器人是什么，并总结了关于机器人检测的现有文献。

\subsection*{2.1 去中心化金融（DeFi）}

以太坊的智能合约支持允许以去中心化的方式执行程序。这催生了涉及加密资产（即代表某种经济资源或价值、基于密码学原语和DLT的数字资产）的复杂功能[3]。去中心化金融（DeFi）使得在像以太坊这样的区块链系统上进行无需信任的借贷或投资成为可能[44]。以下是与本工作相关的DeFi核心组成部分。

代币是通过智能合约创建的加密资产，这些合约遵循特定标准以实现诸如账户间可转移性等核心功能[40]。遵循ERC-20标准的代币是可互换的，即每个单位与同一合约部署的任何其他单位完全相同。根据ERC-721标准实现的智能合约是不可替代代币（NFT）。这意味着它们不可互换，这使得每个NFT都是独一无二的，允许它们被用作数字资产所有权的证明。代币已被用于在区块链上表示美元等价物（稳定币），而NFT已被用于表示艺术品、门票或视频游戏中的角色[41]。

以太坊上的去中心化交易所（DEX）是创建用于相互交换加密资产市场的智能合约。它们的链下对应物，即中心化交易所（CEX），通常用作法定货币与加密资产之间的接口，并通常通过订单簿实现交易。另一方面，DEX通常采用自动化做市商（AMM），允许在没有订单簿的情况下进行交易[25]。为了使这些AMM工作，流行的机制需要流动性来实现交易，并且用户被财务激励以提供流动性[20]。

DeFi中的一个相关概念是最大可提取价值（MEV），最初在Daian等人的开创性工作中被称为矿工可提取价值[13]。它指的是在标准区块奖励和交易费用之外，可以从区块生产中提取的利润。这是通过在一个区块内包含、排除或重新排序交易来实现的[37]。MEV最突出的例子是套利、三明治攻击和清算。套利利用DEX上的价格差异来产生利润。三明治攻击发送一笔交易以推高价格，让受害者在高价买入，然后发送另一笔交易以获利卖出之前购买的东西。清算涉及出售抵押资产以有利的利率偿还债务给买方[28]。

\subsection*{2.2 机器人}

以太坊生态系统中的自动化潜力以及DeFi中涉及的资产数量吸引了被归类为机器人的各种参与者。在像以太坊这样的基于账户的区块链系统背景下，Zwang等人[49]将机器人定义为软件机器人。在这项工作中，我们通过定义互斥的“人类”和“机器人”两类来更精确地描述它们，并建立以下定义2.1：

定义 2.1（机器人）。机器人是一个EOA，它已发送了至少一笔由软件程序编译和发送且没有人类监督的交易，即一笔机器人交易。

\subsection*{2.3 机器人检测与相关文献}

像MEV-inspect和Eigenphi这样的软件程序可以识别与MEV相关的交易。由于这些交易对时间高度敏感，机器人被用来执行它们；因此，我们可以通过识别检测到交易的发送者，将这些程序解释为机器人检测系统。目前，这些系统主要以基于规则的方式关注交易模式，并专注于其代码中建模的一组特定交易。除了这些程序之外，还进行了与识别以太坊上的机器人相关的研究。研究利用启发式方法，允许在无需昂贵注释工作的情况下分析机器人。基于图的启发式方法揭示了机器人类型和私人交易在MEV利用中的普遍性[27]。基于时间的启发式方法在识别自动化钱包行为方面也发挥了重要作用，揭示了交易时序的模式[49]。此外，还使用基于交易和图论的指标对以太坊钱包进行了分类[7]。对参与狙击（MEV的一个子集，旨在尽快获取加密资产以从其他用户推高价格中获利）的机器人进行的分析，通过对交易时序的启发式检查，突显了它们的操作策略和普遍性[9]。总的来说，这些基于启发式的研究为机器人检测提供了坚实的基础，并激发了可在更灵活的基于ML的系统中实施的特征。

MEV机器人因其对以太坊安全和经济的影响而受到特别关注，正如最近的研究所强调的那样[16, 29]。由于时间对利用MEV机会至关重要，机器人被用来自动发现它们并组装和发送针对这些机会采取行动的交易。关键研究已经展示了MEV带来的实际威胁，重点关注实时待处理交易数据分析[13]。存在关于区块链抢先交易的知识系统化，与传统金融相类比，并根据对抗意图对抢先交易进行分类[15]。

虽然我们不知道有任何专门使用ML进行机器人检测的工作，但有几项工作使用ML方法进行更广泛的行为钱包分类。无监督学习，特别是结合随机森林模型的期望最大化算法，已被应用于聚类以太坊钱包和检测异常[4]。关键统计指标已在按以太坊地址分组的数据上计算，并有效地用于基于机器学习的账户分类[4, 17]。有监督和无监督技术已被结合来分类攻击实例，利用深度神经网络进行攻击检测[30]。以太坊上的庞氏骗局检测，结合了手动标注和特征提取，是机器学习在识别欺诈活动中应用的一个例子[11, 17]。此外，比特币交易中的异常检测，采用如一类支持向量机和k-均值聚类等方法，已被综述[34]。

总之，我们发现了适用于机器人检测的软件、采用基于规则的启发式方法检测机器人的研究，以及更普遍地，使用机器学习检测账户特定行为的研究。然而，据我们所知，目前还没有关于以太坊上基于机器学习的机器人检测的研究。这构成了我们进行本研究的理由。

\subsection*{3 机器人类别}

为了全面展示以太坊中DeFi机器人的景观，我们利用从GitHub仓库、示例交易或示例账户收集的现有金融机器人的轶事证据，补充了现有的科学文献（详见附录表5）。我们在图2中报告了我们的分类法，基于收集的证据分为7个大类和24个子类。下面我们概述了所确定的七个类别的特征。

MEV机器人通过快速检测和执行盈利机会来利用区块链状态获利，引发了对公平性和协议完整性的担忧。中心化交易所（CEX）机器人处理中心化交易所与区块链之间的交互。它们包括用于日常运营的热钱包、管理用户资金的存款钱包，以及提供Gas费用覆盖流动性的资金钱包。去中心化交易所（DEX）机器人参与定制交易策略和流动性复投。它们使用自动化机制来交易代币或向池子提供流动性，通常是为了质押奖励。不可替代代币（NFT）机器人参与交易策略和铸造实践。这包括基于定价模型的自动交易和并行铸造，以规避智能合约为防止利用而施加的批量限制。Play-to-Earn（P2E）机器人自动化基于区块链的游戏中用于获取奖励的游戏进程动作。常见动作包括任务、战斗、制作和市场交互。自动化流程使得

\subsection*{4 数据与方法}

在前面的部分中，我们已经定义了机器人并概述了机器人景观。现在，我们描述如何处理公开的区块链数据和智能合约信息以创建我们的数据集。

\subsection*{4.1 数据来源}

我们从以下来源收集数据：

以太坊区块链：我们使用Graphene-lib从Erigon归档节点获取交易数据[18]。我们调查的区块范围是从15,500,000（2022年9月9日）到15,599,999（2022年9月24日），即100,000个区块。这个有限的区间允许我们的方法在消费级硬件上运行。此外，该区间的定义使得它不包含可能使分析产生偏差的显著炒作阶段，并且我们所知的DeFi所有主要功能都已建立，这对许多机器人类别都很重要。虽然此处描述的方法使用了100,000个区块，但区块范围可以任意扩展或限制，以便分别在特征中包含更多信息或允许更快的处理。

Ethereface.io：我们使用Ethereum.io的智能合约函数和事件描述来解码智能合约函数和事件，并使得交换和代币转账中的数据可访问（见附录表7）。

Etherscan.io：使用Etherscan.io的API，我们提供关于已标注地址的额外交易数据，以使标注者能够进行更全面的分析。

\subsection*{4.2 真实数据创建}

我们将区块范围的最后两个区块称为测试区块，并对它们进行标注以生成一个带标签的数据集。为了手动标注在测试区块中发送交易的所有270个EOA，我们雇佣了两位独立的标注者A和B。基于网络文档、索引服务以及自定义聚合指标和图上的信息，两位标注者为所有EOA附加了“机器人”或“人类”标签之一。除了二元标签外，标注者A还添加了细粒度标签，指示机器人的具体用途。如果出现分歧，第三位标注者C根据标注者A和B记录的描述进行最终裁决。在A选择“人类”（因此不提供细粒度机器人标签）而C选择“机器人”的裁决情况下，则B提供细粒度标签。结果得到一个包含137个机器人和133个人类地址的数据集，其中31个地址需要裁决。前两位标注者的Cohen's Kappa为0.77，表明有实质性的共识。最后，为了能够分析区分多个机器人类别的分类器，我们收集了关于长尾MEV机器人的额外数据。我们在整个区块范围上运行MEV-inspect，为套利、三明治和清算机器人每个类别获取111个机器人。每个类别的机器人数量在设计上是相同的，并受限于清算类别中找到的最少机器人数量。

\subsection*{4.3 特征工程}

我们通过将地址的交易历史转换为聚合指标来执行特征工程。这主要是针对传出交易或代币转账，在某些情况下也针对传入交易。我们为每个钱包计算了总共83个特征。这些指标在某些情况下借鉴了交易机器人或人类行为的先验知识，而在其他情况下，它们反映了地址的整体活动，为后续模型提供高度信息丰富的特征。附录中的图6显示了特征的概览，这些特征分为四个主要类别：基于地址的特征（基于地址的40个字符十六进制表示，不需要交易历史）、基于交易的特征（基于交易数据）、基于函数调用的特征（基于发送给智能合约的交易中的输入数据）以及基于事件的特征（基于智能合约产生的数据）。总体策略是首先从账户的交易历史中提取特定数据，然后应用聚合函数生成一个或多个特征。我们首先介绍多次使用的此类函数，然后描述特征，这些特征可能需要直接与特征一起定义的特定函数。

通用函数。函数BenfordsLaw应用于一组给定值，利用本福特定律，该定律预测自然数据集中第一位数字的分布。该函数应用皮尔逊卡方检验来评估第一位数字分布与本福特定律的吻合程度，生成一个p值作为特征[6, 12]。为了捕捉人类交易者使用整数作为认知参考点的倾向，我们定义TradeValueClustering为分组数据中整数和非整数的比率[12]。如果一个数字在第7位有效数字之后的数字为零，我们称其为整数。GapBasedSleepiness计算一个时间戳列表的时间差。时间戳被划分为两天间隔以最小化时区影响并考虑人类睡眠模式。在每个间隔中，确定最大时间差，然后在所有间隔上取平均值以产生一个单一值。Categorical接受一组分类值，并返回类别分布熵，每个类别的发生百分比以及众数。Numerical接受一组数值，并返回平均值、众数、标准差、最小值、最大值和95\%分位数。

\subsubsection*{4.3.1 基于地址的特征。} 特征NLeadingZeros计算以太坊十六进制地址中的前导零数量。具有更多前导零的地址（通常为了效率而挖掘）表明有意创建以减少Gas费用并提高智能合约效率[1]。DigitEntropy测量以太坊地址中数字的熵，捕捉随机性或模式，这些在机器人和人类中可能表现不同，因此具有预测能力。

\subsubsection*{4.3.2 基于交易的特征。} TxPerBlock计算发送的交易数除以总区块数（100,000）。StatisticsTxPerActiveBlock计算每个区块发送的交易数量，移除零值，并应用Numerical。得到的指标反映了活动的强度及其在发送了交易的区块中的分布。

对于基于时间的特征，GapBasedSleepiness分别应用于传入和传出交易的时间戳。TransactionFrequency分别计算传入和传出交易的总数除以第一次和最后一次交易之间的秒数。StatisticsTimeDifference表示Numerical函数应用于连续发送交易之间时间差组后返回的特征。EntropyTime分别对传入和传出交易进行计算，并计算时间戳的熵，在将它们转换为一天中发送的小时后[21]。对于基于价值的特征，BenfordsLaw和TradeValueClustering是针对传出交易中转移的ETH价值计算的。标准字段是分类的“类型”和“状态”，以及数值的“价值”、“gasLimit”和“gasPrice”，连同通过区块中交易数量归一化的交易索引，称为“indexRelative”。对于每个标准字段的数据组，根据数据类型应用Categorical或Numerical。

\subsubsection*{4.3.3 基于函数调用的特征。} 为了过滤与DEX相关的特定智能合约函数调用或发出事件的区块链数据，我们基于Uniswap协议对其进行建模，因为Uniswap是DEX领域的先驱，拥有开源代码库。这导致许多其他协议复制Uniswap的代码和接口，使我们能够仅通过建模Uniswap的函数和事件来捕获它们的功能（见表7）。

我们筛选传出交易历史，查找我们已建模合约的交换。交换对金融机器人尤其重要，因为它们代表了它们使用的许多基本功能，因此可以帮助以特征的形式更好地表示账户行为。我们计算交换中“amountIn”或“amountOut”字段的BenfordsLaw和TradeValueClustering，因为在人类使用DEX网站的情况下，这些字段由用户提供，而智能合约调用中的其他字段通常在交易发送前由网站自动计算。自动计算的值不太适合像TradeValueClustering这样基于认知规则的指标。StatisticsPathLength对所有传出交易获取路径参数，该参数确定输入资产是直接交换为输出资产，还是中间有其他交换，取其长度（直接交换为2），并应用Numerical函数。最后，SwapsPerBlock计算交换交易数除以总区块数。

\subsubsection*{4.3.4 基于事件的特征。} 对于基于UniswapV2的交换，我们通过将“amount0In”和“amount1In”字段相加得到输入金额；对于UniswapV3，我们通过取“amount0”和“amount1”字段的最大值得到输入金额。我们按交换的接收者对数据进行分组，而不是事件数据中的发送者，因为发送者通常是智能合约。我们计算交换中输入金额的BenfordsLaw和TradeValueClustering。此外，在SwapsPerBlock中，给出了交换事件数除以总区块数。类似地，对于具有标准ERC-20转账事件签名的转账事件的价值，计算BenfordsLaw和TradeValueClustering。最后，TransferEventsPerBlock给出转账事件数除以总区块数。

对于基于交换的特征，我们执行降维步骤，将所有建模的函数和事件上的各自特征取平均，将128个特征减少到14个。StatisticsPathLength特征仅对建模的函数取平均。

\subsection*{4.4 数据集创建}

最后，我们使用第4.1节中描述的100,000个区块来创建三个不同大小的数据集，每行代表一个地址及其特征集。

(1) 二分类机器人数据集：我们收集在测试区块中发送交易的所有270个地址来生成此数据集。如第4.2节所述，这些地址与二元标签“机器人”和“人类”相关联。
(2) 聚类数据集：我们收集在100,000个区块窗口内发送交易的所有地址，并排除在测试区块中发送的地址。此数据集是无标签的。
(3) 多类MEV数据集：我们从测试区块中随机采样111个不涉及MEV的地址，并将它们与来自自动标注地址的每个短尾MEV类别的111个地址组合。这产生一个包含4个类别的带标签数据集，每个类别大小为111，对应标签“套利”、“三明治”、“清算”和“非MEV”。

\subsection*{5 结果}

在本节中，我们分析聚类和分类在检测以太坊上机器人方面的有效性。此外，我们研究了在性能最佳的分类器的决策中哪些特征最具影响力。

\subsection*{5.1 聚类}

我们首先使用聚类技术来检测以太坊上的机器人，并探索在无监督环境下机器人和人类的分组效果如何。为了评估聚类的质量，我们测量了所有个体簇中关于机器人和人类两个类别的熵和纯度（见附录A）。纯度是同质性的度量，较高的值意味着该簇更由单一类别的数据点主导。熵是簇中无序程度的度量，对于所有类别的平等表示为1，对于仅包含一个类别的簇为0。因此，对于此分析，我们首先使用无标签的聚类数据集进行拟合，然后使用二分类机器人数据集计算纯度和熵。

在调查了多种聚类方法后，选择了k-均值和高斯混合模型（GMM），因为它们是归纳式的，因此可以用于对二分类机器人数据集中的新实例进行聚类，并且运行速度足够快。对于这两种算法，我们使用scikit-learn实现，其中必须提供簇的数量。我们尝试了不同的方法来确定最佳簇数，例如贝叶斯信息准则的最小化和肘部法，以及轮廓系数。虽然对于k-均值，最小化BIC产生超过100个簇，而使用肘部法在BIC或轮廓系数上产生5到10个簇，但最小化GMM的BIC产生10到30个簇。常见方法建议的这个广泛簇数范围与我们使用的评估指标存在冲突。由于我们有标注数据，我们可以计算关于簇内机器人和人类地址分布的纯度和熵。这些是有吸引力的指标，因为它们表明某个簇是否可以用于选择更高百分比的机器人或人类。冲突在于，更高的簇数，在其他条件相同的情况下，会导致更高的纯度，因为更多的簇平均每个簇的地址更少，这反过来又导致簇内分布纯粹是偶然的机会更大。当一个簇可能仅围绕一个地址时，这种效应变得明显，导致纯度为100\%。因此，我们在5、15和30三个固定簇大小下评估算法，以获得不同模型复杂度水平下的结果。基于BIC优化和“肘部”法的附加结果报告在附录表6中。

对于GMM和k-均值，我们通过最小-最大缩放对聚类数据集中的特征进行预处理到[0, 1]范围，确保特征之间的同等重要性。我们还研究了UMAP嵌入到二维是否会增强性能，以及用列均值或-1填充缺失值是否会产生更好的结果。在缩放后用-1填充缺失值可能是有利的，因为它将特征空间中有缺失值的地址与没有缺失值的地址明显分开。

最后，为了计算纯度和熵，所有预处理方法、聚类方法和簇数的组合都在聚类数据集上拟合，并通过将簇分配给二分类机器人数据集中的地址进行测试。对于整个聚类，我们将这两个指标定义为各个簇指标的加权平均值，权重为相应簇中样本的比例。

表1显示了不同算法、预处理方法和簇大小的结果。具有30个簇、均值填充和无降维的GMM在熵和纯度方面得分最高；我们称之为顶级聚类算法。总的来说，UMAP嵌入对于少量簇是有益的，而对于大量簇的运行则是有害的。对于较小的模型，UMAP嵌入甚至产生了性能最佳的聚类。这可能是由于嵌入丢失了较大模型可以利用的信息，同时通过简化数据仅使用少数簇来帮助有效分离簇。

表2报告了顶级聚类算法生成的包含至少15个地址的6个簇。簇内的多数类基于二分类机器人数据集的细粒度标签。簇6和19的多数类是“人类”，并显示出高纯度。簇5和13也显示出异常高的纯度，多数类分别是CEX和存款钱包。簇2和14质量较低，其熵接近100\%。结论是，存在易于辨别的细粒度机器人类别，如存款钱包和CEX，并且在30个簇的情况下，可以达到82.6\%的纯度。

\subsection*{5.2 分类}

接下来分析不同监督机器学习算法在区分机器人与人类以及其他类型机器人方面的有效性。为此，我们在二分类机器人数据集上评估二分类算法，并在多类MEV数据集上评估多类分类算法。请注意，后者部分是由MEV-inspect生成的，因此在监督环境下并不能轻易地推广到该系统的能力之外。然而，它允许我们进一步测试我们的方法和特征的实用性。为了有效使用所有带标签的数据，我们的方法偏离了典型的数据科学工作流程。我们不是在训练子集上迭代优化模型，然后在保留集上进行评估，而是选择通常在小型表格数据集上表现良好的方法，并仅在整个带标签的数据上评估一次。这使我们能够使用更多数据来展示结果，但限制了优化。

对于我们的二分类和多类分类器，我们使用了scikit-learn实现的随机森林和AdaBoost，以及XGBoost包中的XGBoost[10]。我们将随机森林的估计器数量设置为400，其余参数保持默认值。我们选择这些方法是因为它们已知在小型表格数据集上表现良好[46]。尽管基于树的方法对尺度不敏感，但我们在预处理步骤中应用标准化和均值填充。

为了衡量分类器的性能，我们采用标准的二分类指标：准确率、召回率、精确率和F1，其中机器人代表正类。对于多类问题，我们使用相应指标的宏平均。我们对二分类机器人数据集上的二分类器进行建模。

表3显示，在二分类机器人数据集上测试随机森林、梯度提升和AdaBoost在二元机器人/人类问题上的效用，这些算法的表现相似，特别是在平均准确率的差异和置信区间的大小方面。

表4展示了所调查方法在四类分类问题（套利/三明治/清算/非MEV）上的性能。随机森林表现最佳，我们称之为顶级分类器。更深入地分析顶级分类器，清算是最容易辨别的，准确率达到93\%，捕获Gas使用的特征使得易于区分，因为清算地址在这些特征上显示出异常高的值（如附录图5所示）。值得注意的是，三明治地址很难检测，其中16\%被误分类为套利，13\%被误分类为非MEV，准确率仅为68\%。此外，17\%的非MEV被错误地分类为三明治机器人。再次，随机森林的表现与梯度提升相似，但AdaBoost在二元机器人/人类问题上的表现明显差于前者。总之，我们的特征不仅适用于区分一般机器人与人类，也适用于区分更具体的机器人类别。虽然多类MEV数据集的生成不允许泛化到MEV-inspect能力之外，但我们获得的见解是，我们的特征对各种EOA都具有描述性。

\subsection*{5.3 特征研究}

为了找出在二元和多类设置下顶级分类器最重要的特征，并衡量最重要特征的影响，我们基于二分类机器人数据集和多类MEV数据集，使用可解释AI进行了进一步分析。

\subsubsection*{5.3.1 SHAP值。} SHAP值是机器学习中用于解释模型预测的方法。它们提供了一种方法，通过保持所有其他特征不变，改变所研究的特定特征，来归因每个特征对特定实例最终预测的贡献。在图3中，为了进一步研究顶级分类器，SHAP值被绘制在二分类机器人数据集的样本上，以显示最具影响力的特征。以下是图中特征的描述。

(1) Out-TX-Entropy：传出交易时间分布的熵。在计算熵之前，时间基于一天中的小时进行计数，从而产生离散分布。较高的值倾向于将模型的预测推向机器人类别。
(2) Out-TX-Per-Block：每个区块传出交易的平均数量。此特征表明地址的高活跃度，并且由于高SHAP值地址具有此特征的高值，它似乎会增加模型的怀疑度。
(3) Out-TX-Frequency：在EOA活跃的时间段内，每个区块传出交易的平均数量。高SHAP值地址同时包含此特征值非常高和非常低的地址。低SHAP值地址始终具有低频率。
(4) TX-GasPrice-Max：传出交易中Gas价格的最大值。支付的最高Gas价格越高，模型越倾向于将地址分类为机器人。
(5) Out-TX-Sleepiness：传出交易的GapBasedSleepiness。值越极端，其对模型预测的影响越大。其表示中颜色的平滑过渡表明，它本身就可以改变模型对某一类别的倾向。与其他特征相比，Out-TX-Sleepiness对于高SHAP值的EOA具有较低的特征值。这意味着Out-TX-Sleepiness中的低数值会将模型的预测转向“机器人”。

图4显示了在多类MEV数据集中，顶级分类器对所有四个类别的平均绝对SHAP值。由于SHAP值是针对每个类别单独定义的，因此以更聚合的视图显示。最具影响力的五个特征是：

\subsection*{6 讨论与结论}

我们的研究揭示了金融机器人在以太坊生态系统中的作用。首先，我们识别并分类了在以太坊上活跃的机器人类型，并发布了一个包含133个人类和137个机器人地址的真实数据集。其次，通过部署有监督和无监督的机器学习模型，我们设计了一种自动检测机器人的方法。我们表现最佳的聚类和分类模型分别达到了82.6\%以上的纯度和83\%的准确率。第三，我们使用可解释AI来研究哪些特征对我们的模型最具信息量，发现最有影响力的特征与传出交易的时间、频率、Gas价格和Gas限制有关。

我们注意到，使用高斯混合模型进行聚类在将某些机器人子类型和人类以高纯度分组方面是有效的。这意味着在初始聚类步骤之后，可以将簇映射到某些类别，以筛选出感兴趣的类别。随机森林和梯度提升在二元机器人/人类任务以及分类三种MEV和非MEV的四类任务上都表现良好。分析产生最佳结果的模型所使用的特征，我们发现机器人与人类相比表现出不同的活动/非活动模式，正如本研究首次提出的GapBasedSleepiness特征所捕捉到的那样。

此外，我们的发现有一些相关的影响。首先，我们的分类法建立了一个对现有金融机器人的最先进描述，因此有助于区分那些可以为DLT生态系统提供关键基础设施的机器人和那些可能被用于掠夺性交易和市场操纵的机器人。

其次，尽管我们承认我们考虑的样本很小，但我们的研究为以太坊生态系统中机器人部署的程度提供了初步见解。事实上，我们发现超过50\%的标注地址是机器人。我们的工作可以为这一现象的规模提供额外的见解。通过设计一种自动检测机器人的方法，我们也能够以良好的准确度预测在未来进入的区块中有多少机器人执行交易。类似地，原则上，我们的方法可以扩展到包括以前的交易，从而能够估计这一现象在整个以太坊交易历史中的规模。第三，像SHAP值这样的可解释AI工具可以为人类用户和政策制定者提供有关潜在风险以及可能为机器人的用户特征的信息。

除此之外，我们承认我们的工作面临一些局限性，并为进一步的研究开辟了方向。一个概念性问题是很难准确定义机器人是什么，以及根据其活动精确定义机器人需要哪些条件。然而，我们研究的主要限制是带标签数据集的规模较小。我们目前的真实数据集仅由270个EOA组成，这些EOA由三位独立标注者标注。对此研究的一个直接改进是扩展测试区块和标注者的数量，以增加带标签的数据集并提高所分配标签的置信水平。一个有效的策略可能是采用半监督学习方法，将我们少量的带标签数据与容易获得的未标签数据结合起来。由于局限于一小部分区块会导致高度活跃地址的代表性过高，因此扩展标注的区块数量也将产生更具代表性的分布。在未来的工作中，额外的数据可能允许采用标准的数据科学工作流程，如CRISP-DM，并实现模型选择和参数的优化。此外，我们目前仅分析了100,000个区块。然而，我们的方法很容易扩展到更大的时间窗口，本文的主要目的是提供识别机器人的方法论方法。另一个局限性是本研究关注于从事金融活动的机器人。然而，机器人很可能也在其他领域得到利用，例如，在去中心化自治组织中自动化投票过程。

未来的研究可以将我们的分类法扩展到非金融机器人。此外，我们目前对一类受限的机器人（即MEV机器人）进行多类分类。未来的工作也可以将此分析扩展到更多的机器人类别。最后，未来的研究可以调查其他区块链上的这种现象，并验证不同DLT之间的机器人类别是否存在差异。


参考~\cite{zjuthesisrules, zjuthesis}

\begingroup
    \linespreadsingle{}
    \printbibliography[title={外文翻译参考文献}]
\endgroup
